\documentclass{standalone}
\usepackage{aided_preamble}
\begin{document}

The (static) Electron Density (ED) of a molecule with M positively charged nuclei  and N negatively charged
electrons gives the probability of finding an electron in the infinitesimally small space
$d \mathbf{r_x}$ at the point $\mathbf{r_x} = (x, y, z)$:

$$
\rho(r_x) = N \int \Psi(r_x, \mathbf{r}_2, \ldots, \mathbf{r}_N; \mathbf{R}) \Psi^*(r_x, \mathbf{r}_2, \ldots, \mathbf{r}_N; \mathbf{R}) \, d\mathbf{r}_2 \ldots d\mathbf{r}_N
$$

where $\Psi$ is the spinless ground-state 'wave function' for the N electrons, distributed in the field of nuclei with fixed positions
(given by the 3M vectors $\mathbf{R}$).


Given an equlibrium geometry $\mathbf{R}$ and accompanying ground-state wave function $\Psi$, a common practice
is to describe the many electron wavefunction, $\Psi_i$, as an antisymmetric linear combination of one-electron 
Molecular Orbitals (MOs), $\Psi_i$,
each of which is described by a linear combination of Atomic Orbitals (AO), $\phi_j$.

% ------------ Basis definitions -------------------------------------------
An atomic orbital (AO) is a contraction of primitive Gaussians:
\begin{equation}
\label{eq:phi_j}
\phi_j(\mathbf r)=\sum_{k} g_{jk}\,g_k(\mathbf r)
\end{equation}
Molecular orbitals (MOs) are written directly in the primitive basis
using the coefficients that appear in the code:
\begin{equation}
\label{eq:Psi_i}
\Psi_i(\mathbf r)=\sum_{k} C_{ik}\,g_k(\mathbf r)
\end{equation}

% ------------ Density in the AO basis -------------------------------------
The electron density is
\begin{equation}
\label{eq:rho_phi}
\rho(\mathbf r)=\sum_{\mu,\nu}P_{\mu\nu}\,
                \phi_\mu(\mathbf r)\,\phi_\nu(\mathbf r),
\qquad
P_{\mu\nu}=2\sum_i n_i\,c_{i\mu}\,c_{i\nu}\
\end{equation}
(AO-level MO coefficients \(c_{i\mu}\) are related to \(C_{ik}\) by
\(C_{ik}= \sum_\mu c_{i\mu}\,g_{\mu k}\).)

% ------------ Transform to the primitive basis ----------------------------
In \textsc{AIDED} we insert \eqref{eq:phi_j} into \eqref{eq:rho_phi}:
\begin{align}
\rho(\mathbf r)
  &=\sum_{\mu,\nu}P_{\mu\nu}
    \!\Bigl(\sum_{k} g_{\mu k}\,g_k(\mathbf r)\Bigr)
    \!\Bigl(\sum_{l} g_{\nu l}\,g_l(\mathbf r)\Bigr)\notag\\
  &=\sum_{k,l}
    \Bigl(\sum_{\mu,\nu}g_{\mu k}\,P_{\mu\nu}\,g_{\nu l}\Bigr)
     g_k(\mathbf r)\, g_l(\mathbf r).
\end{align}

Define the primitive-basis density matrix
\begin{equation}
\label{eq:Dkl}
D_{kl}=\sum_{\mu,\nu}g_{\mu k}\,P_{\mu\nu}\,g_{\nu l},
\end{equation}
so that
\[
  \rho(\mathbf r)=\sum_{k,l}D_{kl}\,g_k(\mathbf r)\,g_l(\mathbf r)
  =g^{\mathsf T}D\,g.
\]

% ------------ One-shot construction used in the code ----------------------
Because Gaussian’s *.wfn file already folds the contraction coefficients
into the exported MO matrix \(C_{ik}\), we may build \(D\) in a \emph{single
pass} exactly as in the Python routine:

\begin{equation}
\label{eq:denmat}
D_{kl}= \sum_{i} n_i\,C_{ik}\,C_{il}.
\end{equation}

The entire grid calculation then reduces to evaluating
\(\rho(\mathbf r)=g^{\mathsf T}Dg\) with a \textbf{fixed} \(D\), giving the
computational speed-up described in the text.

\end{document}
